% Slajd 5 – źródła danych, aktualność, wiarygodność (kryterium: wiarygodność oraz sposób prezentacji i aktualizacji danych, 15 %).
\begin{frame}{Dane, którym można zaufać}
  \begin{columns}[T,onlytextwidth]
    \begin{column}{0.49\textwidth}
      \begin{block}{3 źródła informacji}
        \begin{itemize}
          \setlength{\itemsep}{0.3mm}
          \item \textbf{Open Data \& OSM}\\
                Dane publiczne i~otwarte
          \item \textbf{Property Owner}\\
                Potwierdzone informacje bezpośrednio od właściciela obiektu
          \item \textbf{Community}\\
                Użytkownicy potwierdzają, poprawiają i~uzupełniają dane
        \end{itemize}
      \end{block}
    \end{column}
    \begin{column}{0.49\textwidth}
      \begin{block}{Każdy fakt ma kontekst}
        \centering\textbf{Źródło \textperiodcentered\ Data \textperiodcentered\ Status weryfikacji}\par
      \end{block}
      \vspace{1mm}
      \begin{exampleblock}{Accessly Trust Model}
        \pill{accNavy}{Owner verified} $\rightarrow$ \pill{accOk}{Community confirmed} $\rightarrow$\\[0.8mm]
        \pill{accNavy}{Official source} $\rightarrow$ \pill{accMuted}{Unverified}
        \begin{itemize}
          \setlength{\itemsep}{0pt}
          \item nieaktualne dane są oznaczane,
          \item sprzeczności są widoczne,
          \item brak informacji nigdy nie oznacza „dostępne”.
        \end{itemize}
      \end{exampleblock}
    \end{column}
  \end{columns}

  \note[item]{Przypadek testowy wymagany w~wyzwaniu – dane sprzeczne: właściciel mówi „bez schodów”, dwie osoby zgłaszają „nieaktualne” $\rightarrow$ status „Wymaga ponownej weryfikacji”, wartość zostaje z~ostrzeżeniem, obie strony widać w~historii.}
  \note[item]{Przypadek – źródło niedostępne: usługa ArcGIS miasta nie odpowiada $\rightarrow$ aplikacja pokazuje warstwę ze snapshotu z~datą „stan na”; żaden ekran nie jest pusty.}
  \note[item]{Przypadek – dane niepełne: „Spełnia 1~z~3 Twoich potrzeb, brak danych: szerokie drzwi” – miejsce nie dostaje odznaki „Dopasowane do Ciebie”.}
  \note[item]{Sprawdziliśmy też GTFS Krakowa i~dane.gov.pl: pola o~dostępności są puste, nikt nie publikuje awarii wind – dlatego awarie pochodzą ze zgłoszeń.}
\end{frame}
